\par\noindent\begin{minipage}{\linewidth}
% !TeX root = ../教材题库_全部条目_转写工作稿.tex
% Source: XJ-B1-C02-PR-P042-Y1011
% 第2章 / 2.1 相等关系与不等关系 / 教材 p.42 / PDF 第49页
% 已对照原页转写题干及公式；待教师复核；未制作答案
% 原题及相邻原解截图（核对用）：../题目原图/XJ-B1-C02-PR-P042-Y1011-01.webp
% 本文件只在转写完成后作为题干使用；OCR 条目尚不是可靠题干。
\begin{bankitem}{p.42 练习 第2题}{第2章 / 2.1 相等关系与不等关系 / 教材 p.42 / PDF 第49页}{XJ-B1-C02-PR-P042-Y1011}{已对照原页转写题干及公式；待教师复核；未制作答案}
如图，动物园要以墙体为背面，用钢筋网围成四间面积相同的矩形虎笼。
\bankfig{C02-PR042}{0.58}
\subq{1}现有可围 $36\,\mathrm m$ 长钢筋网的材料，每间虎笼的长、宽各设计为多少时，可使每间虎笼的面积最大？
\subq{2}若每间虎笼的面积为 $20\,\mathrm{m^2}$，则每间虎笼的长、宽各设计为多少时，可使围成四间虎笼的钢筋网总长最小？
\end{bankitem}
\end{minipage}\par\medskip
